% sparc_shells_body.tex
% =====================
% Journal-neutral manuscript body. Included from:
%   sparc_shells.tex  (article-class standalone)
%   aj_wrapper.tex    (aastex701)
%   mnras_wrapper.tex (mnras)
% Do not include \documentclass, \begin{document}, or \end{document}.
% Do not include preamble \usepackage commands; those live in the wrappers.

\title{Localized Residual Structure in SPARC Rotation Curves: A Burkert-Backbone Model with BIC-Selected Shells}

\author[orcid=0009-0004-1153-2464]{Ronald Bibb}
\affiliation{Independent Researcher, Lilburn, GA, USA}
\email[show]{ronbibb@gmail.com}

\begin{abstract}
We test whether localized residual structure in galaxy rotation curves can improve halo modeling beyond conventional smooth profiles. Using 102 SPARC galaxies (T=2--9), we implement a framework consisting of a Burkert halo plus zero, one, or two Gaussian shells selected by Bayesian Information Criterion (BIC), with shell width constrained by $\sigma/r \leq 0.4$ to prevent shells from mimicking an additional smooth halo component.

The framework adequately fits 91/102 galaxies, compared with 65/102 for Burkert-only and 52/102 for free-concentration NFW. Shell-bearing fraction shows a significant morphology dependence (per-galaxy T-label permutation $p_{\rm two\mbox{-}sided}=0.002$; per-bin Spearman $\rho=-0.833$, $p=0.010$; Mann-Kendall $p=0.003$), with a strong T=2 excess and lower shell fractions at late types. Synthetic null tests show that smooth-Burkert and smooth-NFW truth cannot reproduce the observed pattern: Burkert-truth false-positive rates remain uniformly low and fall far below real-data shell-bearing fractions in every T-bin, while NFW-truth produces high but non-monotonic shell-selection rates inconsistent with the observed trend.

Replacing the Burkert backbone with the more flexible Einasto profile preserves both primary findings: the morphology gradient and the requirement for localized structure in NGC~5055, where Burkert, NFW, DC14, and Einasto smooth profiles all fail to reproduce the observed rotation-curve structure. The morphology gradient also survives T-dependent stellar mass-to-light variation, refuting stellar-population systematics as a primary driver.

We interpret the framework phenomenologically. The results indicate that a significant subset of SPARC rotation curves contains localized residual structure not readily absorbed by smooth halo parameterizations alone, though the physical origin of these features remains unresolved.
\end{abstract}

\keywords{Galaxy rotation curves (619) --- Galaxy dark matter halos (1880) --- Spiral galaxies (1560) --- Disk galaxies (391) --- Galaxy kinematics (602)}

% =============================================================================
\section{Introduction}
\label{sec:intro}

Galactic rotation curves provide one of the clearest empirical signatures of mass discrepancies in disk galaxies. Across a wide range of galaxy masses and morphologies, observed circular velocities remain elevated beyond the radius where luminous baryons alone would predict decline. This behavior is conventionally modeled by embedding galaxies in smooth dark matter halos, with commonly used profiles including NFW \citep{NFW1996,NFW1997}, Burkert \citep{Burkert1995}, pseudo-isothermal, Einasto, and related variants. Tight relationships between baryonic and dynamical properties have also been characterized statistically across SPARC \citep{Lelli2016,McGaugh2016}, providing a useful population-level reference against which residual structure can be assessed.

Smooth profiles are powerful because they capture the broad radial structure of dark matter halos with few parameters. However, many high-quality rotation curves also exhibit localized deviations from smooth behavior: inner excesses, mid-disk inflections, shoulders, and outer flattening features (for a review of rotation-curve diversity and core/cusp systematics, see \citealt{deBlok2010}; observational evidence for cores in late-type dwarfs is reviewed in \citealt{Oh2015}). These structures are often absorbed into global halo parameters, attributed to baryonic uncertainty, or treated as object-specific irregularities. Yet if such features occur systematically, then smooth profiles may be missing an architectural component of halo structure rather than merely requiring better parameter tuning.

In this work, we test the hypothesis that some rotation-curve features reflect localized residual structure that smooth halo profiles cannot capture. We do not assume that individual features can be uniquely associated with specific historical events, nor that they are necessarily of halo origin. Instead, we ask a population-level question: does allowing a small number of localized residual components improve rotation-curve fits in a statistically preferred way, and does the occurrence of such components correlate with galaxy morphology in a manner consistent with hierarchical assembly?

We construct a minimal architectural extension of a smooth halo model. The base profile is a Burkert halo, chosen because it provides a flexible cored backbone for disk-galaxy rotation curves. We then allow zero, one, or two Gaussian shell components, with shell number selected per galaxy by the Bayesian Information Criterion. The shell width is constrained by $\sigma/r \leq 0.4$ so that shells remain localized and cannot simply mimic an additional smooth halo component. When zero shells are selected, the model reduces exactly to Burkert-only. Throughout this paper we use ``shell'' as a generic label for these BIC-selected localized Gaussian residual components in the model. The terminology is not intended to invoke a specific physical origin (tidal-debris stellar shells in the sense of \citealt{HelmiWhite1999}, accretion remnants, or otherwise); the components are mass-density features in the dark matter inferred from rotation-curve residuals, and their physical interpretation is deferred to \S\ref{sec:discussion}.

We apply this framework to 102 SPARC \citep{Lelli2016} galaxies with morphological types T=2--9, spanning Sab through Sdm spirals. The comparison models are Burkert-only and free-concentration NFW, each fit on the same data subsets with the same multi-restart optimization procedure.

The central result is that the Burkert-plus-shell framework adequately fits 91/102 galaxies, compared with 65/102 for Burkert-only and 52/102 for NFW. More importantly, shell-bearing fraction shows a statistically significant declining trend with morphological type (per-galaxy T-label permutation test on 102 individual galaxies, $p_{\rm two\mbox{-}sided} = 0.002$; supporting per-bin tests on 8 morphology bins give Spearman $\rho = -0.83$ and Mann-Kendall $p = 0.003$, with bootstrap 95\% CI on $\rho$ from per-galaxy resampling $[-0.95, -0.22]$), driven primarily by a statistically significant excess in early-type (T=2) galaxies and a low fraction at the latest types. This morphology gradient is the primary empirical content of the paper. We support its statistical robustness with a synthetic null test (\S\ref{sec:null}) showing that the Burkert-truth false-positive rate is uniformly low ($\leq 5\%$ in 6 of 8 T-bins), an order of magnitude below the real-data shell-bearing fractions in every T-bin --- ruling out smooth-Burkert truth as the dominant driver of the gradient, with NFW-truth ruled out by analogous reasoning, though mixed-truth and external-family scenarios remain logically possible.

We emphasize that the present framework is phenomenological. It does not prove that any specific fitted shell corresponds to a particular merger event, satellite accretion episode, or dark-sector structure, and it does not establish that the localized components are necessarily of halo origin --- they could equally reflect localized baryonic features not captured by the SPARC mass-to-light decomposition, kinematic artifacts of triaxial halo or warped-disk geometry, or non-circular motions. One-dimensional rotation curves alone do not contain enough information for such per-galaxy attribution. The claim is more limited but also more robust: BIC-selected localized residual components capture rotation-curve features that smooth profiles systematically miss, and the population-level occurrence of those components varies with morphological type in a manner consistent with hierarchical-assembly expectations and inconsistent with the simplest smooth-halo null scenarios.

The framework is also agnostic with respect to the dark-matter-versus-modified-gravity debate. SPARC has been the primary observational anchor for the radial acceleration relation and MOND-style interpretations \citep[][for a review of MOND in the rotation-curve context, see \citealt{FamaeyMcGaugh2012}]{McGaugh2016}, in which rotation-curve regularities are recast as an empirical $g_{\rm obs}(g_{\rm bar})$ relation rather than the mass distribution of an unseen halo. Our analysis treats $V^2_{\rm DM} \equiv V^2_{\rm obs} - V^2_{\rm bar}$ as the residual signal to be modeled and does not attempt to distinguish between dark-matter and modified-gravity interpretations of that residual; either can in principle host the localized residual structure we report. Whether shell-bearing galaxies depart systematically from the radial acceleration relation, or whether shell radii correlate with the radii at which $g_{\rm obs}/g_{\rm bar}$ deviates most strongly from the relation, are questions left for future work.

% =============================================================================
\section{Methods}
\label{sec:methods}

\subsection{Sample and data}
\label{sec:sample}

The SPARC database \citep{Lelli2016} provides Spitzer Photometry \& Accurate Rotation Curves for 175 disk galaxies. We restrict our analysis to morphological types T=2--9, spanning Sab to Sdm spirals, yielding a sample of 102 galaxies with 2334 total rotation curve data points. The original analysis was conducted on T=2--7 (80 galaxies); we extended to T=8--9 (an additional 22 galaxies) to test whether the morphology gradient continues into the late-type regime where merger histories are predicted to be poorest.

For each galaxy, the SPARC rotation files provide observed velocities $V_{\rm obs}(r)$ with uncertainties $\sigma_V(r)$, gas velocity contributions $V_{\rm gas}(r)$, stellar disk contributions $V_{\rm disk}(r)$, and stellar bulge contributions $V_{\rm bulge}(r)$. The squared baryonic velocity is constructed using fixed mass-to-light ratios:
\begin{equation}
V^2_{\rm bar}(r) = V_{\rm gas}|V_{\rm gas}| + 0.5 V_{\rm disk}|V_{\rm disk}| + 0.7 V_{\rm bulge}|V_{\rm bulge}|
\label{eq:vbar}
\end{equation}
with $\Upsilon_{\rm disk} = 0.5$ and $\Upsilon_{\rm bulge} = 0.7$ in solar units at 3.6~$\mu$m. The dark matter velocity contribution is then
\begin{equation}
V^2_{\rm DM}(r) = V^2_{\rm obs}(r) - V^2_{\rm bar}(r).
\label{eq:vdm}
\end{equation}

When $V^2_{\rm bar}(r) \geq V^2_{\rm obs}(r)$, the inferred $V^2_{\rm DM}$ is unphysical. Such points are excluded from the fit.

\subsection{Framework architecture}
\label{sec:arch}

We model the dark matter velocity-squared as the sum of a Burkert backbone and zero or more discrete Gaussian shells:
\begin{equation}
V^2_{\rm DM}(r) = V^2_{\rm Burkert}(r; \rho_0, a) + \sum_i V^2_{{\rm shell},i}(r; M_i, r_i, \sigma_i).
\label{eq:framework}
\end{equation}

The Burkert profile is
\begin{equation}
\rho_{\rm Burkert}(r) = \frac{\rho_0 a^3}{(r + a)(r^2 + a^2)}.
\end{equation}

Each shell is a Gaussian mass distribution with total mass $M_i$, central radius $r_i$, and width $\sigma_i$:
\begin{align}
M_{\rm shell}(r) &= \tfrac{1}{2} M_i \left[ \mathrm{erf}\!\left(\frac{r - r_i}{\sigma_i\sqrt{2}}\right) - \mathrm{erf}\!\left(\frac{-r_i}{\sigma_i\sqrt{2}}\right) \right] \\
V^2_{\rm shell}(r) &= G M_{\rm shell}(r) / r.
\end{align}

We constrain shells to be physically localized through three architectural choices: (1) $\sigma_i / r_i \leq 0.4$, enforced strictly via the reparameterization $\sigma_i = f_i \, r_i$ with $f_i \in [0.01, 0.4]$ as a box-bounded fit parameter; (2) $10^6$~M$_\odot \leq M_i \leq 5 \times 10^{10}$~M$_\odot$; (3) $N_{\rm shells} \in \{0, 1, 2\}$. The radial bound on $r_i$ is set in the multi-restart fit procedure to $r_{\rm max} \in \{3, 6, 12\}$~kpc, where each value is used as a separate starting condition and the resulting fits are compared by BIC. When $N_{\rm shells} = 0$, equation~(\ref{eq:framework}) reduces to $V^2_{\rm DM} = V^2_{\rm Burkert}$. Because the BIC selection (\S\ref{sec:fitting}) can choose zero shells, the framework reduces exactly to Burkert-only whenever shell components are not warranted by the data.

\subsection{Fitting procedure}
\label{sec:fitting}

For each galaxy, fits are performed using non-linear least squares (\texttt{scipy.optimize.curve\_fit}) minimizing
\begin{equation}
\chi^2 = \sum_j \left[\frac{V_{{\rm obs},j} - V_{{\rm model},j}}{\sigma_{V,j}}\right]^2
\end{equation}
with $\sigma_V$ floored at 1~km/s. We adopt a multi-restart procedure with 8--24 distinct initial-condition vectors per fit. The \texttt{maxfev} parameter is set to 15{,}000--20{,}000 (default is 1{,}000), and the fit returning the lowest $\chi^2$ is retained.

For each galaxy, we fit the framework with $N_{\rm shells} = 0, 1, 2$ and select among configurations using the Bayesian Information Criterion:
\begin{equation}
\mathrm{BIC}(N) = \chi^2_{\min}(N) + p(N) \ln(n)
\end{equation}
where $p(N) = 2 + 3N$ and $n$ is the number of data points. We follow \citet{KassRaftery1995} for interpretation of $\Delta\mathrm{BIC}$ differences in cross-model comparisons: $|\Delta\mathrm{BIC}| < 2$ inconclusive, 2--6 positive, 6--10 strong, $>10$ very strong evidence.

\subsection{Comparison models}
\label{sec:comparison}

We fit two smooth-halo models for direct comparison. \textbf{Burkert-only} uses the same parameterization as the backbone, without shells. \textbf{NFW} \citep{NFW1996,NFW1997} uses $\rho_{\rm NFW}(r) = \rho_s / [(r/r_s)(1 + r/r_s)^2]$ with free $\rho_s$ and $r_s$. We do not fix concentration via cosmologically-motivated mass-concentration scalings \citep[e.g.,][]{DuttonMaccio2014}, because removing this degree of freedom disadvantages NFW at small samples.

\subsection{Convergence verification}

Single-restart procedures of the kind often used in nonlinear profile fitting (with \texttt{maxfev} defaults of 1{,}000--2{,}500) systematically converged to suboptimal local minima for galaxies requiring shell components. Under single-restart, NGC~2403 yielded a framework BIC of 728, while multi-restart reduced it to 193. NGC~2841 yielded NFW BIC of 1547 under single-restart, reduced to 160 under multi-restart. All fits reported here use the multi-restart procedure.

% =============================================================================
\section{Results}
\label{sec:results}

\subsection{Aggregate fit quality across T=2--9}

Across 102 galaxies and 2334 data points, the framework reduces total $\chi^2$ by 80.0\% versus Burkert-only and 76.2\% versus NFW (Table~\ref{tab:aggregate}). A first property worth flagging before quoting adequacy counts: BIC selects zero shells in 50 of 102 galaxies. In nearly half the sample, the framework reduces exactly to Burkert-only --- the BIC parameter penalty effectively suppresses shells when the data does not warrant them, and the framework is not biased toward over-fitting localized structure. The shell-bearing classifications discussed below therefore reflect genuine model preference, not an architectural tendency to add shells regardless.

We adopt $\chi^2_{\rm red} < 1.5$ as a threshold for adequate fit. This threshold is intentionally permissive relative to the formal unit-variance expectation: SPARC velocity uncertainties are known to be modestly underestimated, particularly at outer radii where beam smearing and asymmetric drift corrections introduce systematic uncertainty not captured in the quoted $\sigma_V$ \citep[][\S 3.2]{Lelli2016}; $\chi^2_{\rm red} \approx 1$--$1.5$ corresponds in practice to fits that lie within the visible scatter of the data points. We use the same threshold for all four models so that the comparison remains internally consistent. Because adequacy counts are inevitably threshold-dependent, Table~\ref{tab:aggregate} also reports the count at the stricter $\chi^2_{\rm red} < 1.0$ threshold. The framework's dominance is threshold-independent within this range: at the stricter cut, 86/102 framework, 57/102 Burkert, 41/102 NFW; at the headline $\chi^2_{\rm red} < 1.5$ cut, 91/102 framework (89.2\%), 65/102 Burkert-only (63.7\%), and 52/102 NFW (51.0\%). The gap between the framework and the smooth alternatives actually widens slightly under stricter scoring.

The framework rescues 39 galaxies that NFW fails and 26 galaxies that Burkert fails. The reverse --- galaxies where NFW or Burkert adequately fit but the framework fails --- occurs in zero galaxies, the latter as a consequence of the framework's superset architecture. At the late-type extreme (T=9), Burkert clearly outperforms NFW (median NFW/Burkert $\chi^2$ ratio $\approx 4.6$, Burkert wins per-galaxy in 14 of 16 cases); base-profile coupling concerns are minimal at T=9.

\subsection{Win/loss analysis}

By $\Delta\mathrm{BIC}$, the framework is very strongly preferred ($\Delta\mathrm{BIC} < -10$) over NFW in 53 galaxies and over Burkert in 38 galaxies (Table~\ref{tab:dbic}, Fig.~\ref{fig:dbic}). The framework-vs-Burkert comparison is structurally asymmetric: when BIC selects zero shells, the framework reduces exactly to Burkert-only, so a ``loss'' against Burkert can occur only via the BIC penalty for a shell that the optimizer attempted to add and BIC then rejected --- a specific kind of inefficiency rather than a substantive failure. The 55 ``inconclusive'' galaxies in the FW-vs-Burkert column reflect this structural relationship: in 50 of those 55, the framework's BIC-selected configuration is zero-shell (i.e., framework = Burkert exactly).

The framework-vs-NFW comparison is therefore the genuinely informative test of the architecture's contribution. There, 8 of 102 galaxies BIC-prefer NFW over the framework (6 in ``Other positive'' plus 2 in ``Other strong''), though never very strongly --- the ``Other very strong'' bucket is empty. The asymmetry between FW-vs-NFW preference (53 strongly preferring FW vs 8 preferring NFW, none very strongly) is large but not absolute. Median $\Delta$BIC values for the FW-vs-NFW comparison per Table~\ref{tab:dbic} bucket quantify the depth of preference within each verdict: $-34.8$ for ``FW very strong'' (n=53), $-7.7$ for ``FW strong'' (n=6), $-3.2$ for ``FW positive'' (n=16), $-0.7$ for ``Inconclusive'' (n=19), $+3.2$ for ``Other positive'' (n=6), $+8.5$ for ``Other strong'' (n=2). The asymmetry is therefore not just in counts: when the framework wins it typically wins by a large BIC margin, whereas when NFW wins the margin is shallow.

The 8 NFW-preferring galaxies do not form a distinctive subpopulation by morphology, mass, or kinematics. T-type, $V_{\rm flat}$, and $\log M_\star$ distributions in the 8 are consistent with the full sample; notably none are at T=2 (where shell-bearing fraction is highest), and otherwise they scatter across T=3--9. The 8 split into two structurally different subsets. In 4 galaxies (NGC~4157, NGC~3992, UGC~06446, NGC~4183) BIC selects zero shells, so the framework reduces exactly to Burkert; the NFW preference in these cases reflects a Burkert-versus-NFW disagreement on the smooth-profile shape, not a shell-related failure. In the remaining 4 (UGC~02259, UGC~00128, UGC~08490, UGC~12506) the framework selects shells but the BIC parameter penalty tips the comparison toward NFW; in 3 of these 4 the framework's $\chi^2_{\rm red}$ is in fact lower than NFW's, so the NFW preference reflects BIC's parsimony pressure rather than a genuine fit-quality advantage. Seven of the 8 have both NFW and the framework adequate at the stricter $\chi^2_{\rm red} < 1.0$ threshold, indicating these are galaxies for which multiple smooth-profile descriptions fit comparably well and the BIC margin is genuinely soft.

\begin{figure*}
\centering
\includegraphics[width=\textwidth]{figures/figure5_dbic_histogram.pdf}
\caption{$\Delta$BIC distribution across the SPARC T=2--9 sample (n=102). The histograms make visible the asymmetry: zero galaxies fall in the very strong evidence against framework region ($\Delta\mathrm{BIC} > +10$) for either comparison.}
\label{fig:dbic}
\end{figure*}

\subsection{T-type morphology gradient}
\label{sec:morphology}

Shell-bearing fraction varies systematically with morphological type (Table~\ref{tab:morphology}, Fig.~\ref{fig:morphology}). The strongest test of the trend's significance is a permutation test in which T-labels are shuffled across the 102 individual galaxies (10{,}000 resamples, holding shell-bearing flags fixed): observed Spearman $\rho = -0.83$ against a null distribution centered on zero, giving $p_{\rm one\mbox{-}sided} = 0.001$, $p_{\rm two\mbox{-}sided} = 0.002$. This is the most independent of our trend tests because it operates on per-galaxy data rather than per-bin summaries. As supporting evidence, two rank-based monotonic trend tests on the per-bin fractions also reject the no-trend null (Spearman analytic $p = 0.010$; Mann-Kendall $p = 0.003$); these two tests are not independent of each other because both operate on the same eight-bin vector. A non-parametric bootstrap (10{,}000 galaxy-level resamples) gives a 95\% confidence interval on $\rho$ of $[-0.95, -0.22]$, excluding zero.

The shape of the gradient is worth describing precisely. T=2 (Sab) galaxies exhibit shell-bearing fraction of 100\%, while T=9 (Sdm) --- disk-dominated dwarfs --- show 31\%, the lowest in the sample. The intermediate types T=3--6 cluster tightly near 54\% (range 53--56\%), then drop to 36\% at T=7, 33\% at T=8 ($n=6$), and 31\% at T=9. The trend signal is therefore driven primarily by the T=2 excess and the late-type (T=7, T=8, T=9) decline; the middle of the gradient is essentially flat. We use ``morphology gradient'' to refer to this overall pattern, not to claim per-bin monotonicity.

As a robustness check on the gradient's dependence on the T=2 bin ($n=9$), we re-ran the trend tests excluding T=2 (T=3--9, $n=93$, 43 shell-bearing). Per-galaxy permutation gives $p_{\rm two\mbox{-}sided} = 0.094$; per-bin Spearman $\rho = -0.75$, $p = 0.052$; Mann-Kendall $p = 0.069$. The full-sample monotonic-trend tests are therefore dominated by the T=2 excess. The empirical pattern is more accurately described as two contrasts that are independently significant: T=2 (100\%) vs T=3--6 (31/57 = 54\%) gives Fisher one-sided $p = 0.007$, and T=3--6 (54\%) vs T=7--9 (12/36 = 33\%) gives Fisher one-sided $p = 0.038$. The late-type real-data fraction (33\%) also exceeds the Burkert-truth null expectation (1/58 = 1.7\%) at Fisher one-sided $p = 2.5 \times 10^{-5}$. Both pieces of the gradient --- early-type excess and late-type deficit --- are statistically robust independently of each other, even though the full-sample monotonic-trend tests depend on T=2 carrying most of the dynamic range.

\begin{figure*}
\centering
\includegraphics[width=\textwidth]{figures/figure4_shellfrac_vs_T.pdf}
\caption{Shell-bearing fraction in real SPARC data (blue) vs synthetic-null expectation (gray) by morphological type, extended to T=2--9 (n=102 galaxies). Errorbars are 1$\sigma$ binomial counting uncertainties. Fisher one-sided $p$-values (where null comparison was run, T=2--7) are annotated above each bin.}
\label{fig:morphology}
\end{figure*}

Hierarchical assembly theory predicts that early-type galaxies have richer accretion histories than late-type galaxies, plausibly leaving more abundant localized halo features. The framework's shell components --- interpreted phenomenologically as localized corrections to the smooth backbone --- are observed to occur preferentially in early-type galaxies, consistent with this prediction. We do not claim a unique causal link between any specific shell and any specific event; we report the population-level association.

\subsection{Single-galaxy showcase: NGC 5055}
\label{sec:ngc5055}

NGC 5055 (M~63), at distance $D = 9.90$~Mpc, has T-type 4 and asymptotic flat rotation velocity $V_{\rm flat} = 179$~km/s. We use it as a single-galaxy showcase because its rotation curve has been well-studied across multiple decades and its inner kinematic structure provides one of the cleanest tests in this sample of where smooth halos succeed and fail. The case is illustrative rather than load-bearing: NGC~5055 demonstrates at the individual-galaxy level a pattern that the morphology-gradient analysis (\S\ref{sec:morphology}) and the universal-failure subset (\S\ref{sec:universal}) establish at the population level. Our headline statistical claims rest on the 102-galaxy sample, not on this one object.

The rotation curve includes 22 usable data points. Fitting Burkert, NFW, and the framework with the multi-restart procedure (and additionally DC14 as discussed below) yields the values shown in Table~\ref{tab:ngc5055}. The framework's BIC-preferred configuration includes one shell at $r = 12.0$~kpc with $\sigma = 2.91$~kpc and $M = 3.78 \times 10^{10}$~M$_\odot$. The shell sits well within the disk and at the radius where the rotation curve shows its largest residual relative to a smooth Burkert fit. The fitted shell mass is dynamically substantial, indicating that any model component accounting for the residual at this radius must carry a non-negligible mass contribution, regardless of physical interpretation. We do not claim per-galaxy event identification: the localized component could represent a genuine mass concentration of any physical origin, a localized baryonic feature not captured by the SPARC stellar/gas decomposition, or a kinematic artifact of triaxial halo or warped-disk geometry that 1D azimuthal averaging cannot disentangle from intrinsic mass distribution. The $\chi^2_{\rm red}$ comparison establishes that smooth halos (Burkert, NFW, DC14) cannot adequately fit this rotation curve; it does not establish what physical structure the localized component actually is.

To test whether a more sophisticated cosmologically-motivated smooth profile can capture NGC 5055's structure, we additionally fit the DC14 \citep{DiCintio2014a,DiCintio2014b} feedback-modified profile, which transforms cusps into cores at intermediate stellar mass ratios. DC14 with three free parameters ($\rho_s$, $r_s$, $X = \log_{10}(M_\star/M_{\rm halo})$) achieves $\chi^2_{\rm red} = 9.01$, essentially identical to Burkert's 9.53. The optimizer converged on a strongly cored profile ($\gamma = 0.24$) with $X = -2.54$, well within the DC14-valid range, confirming that the failure is not due to NFW-vs-Burkert shape choice. The framework with one BIC-selected shell at the same radius achieves $\chi^2_{\rm red} = 1.46$.

\begin{figure}
\centering
\includegraphics[width=\columnwidth]{figures/figure2_ngc5055_showcase.png}
\caption{NGC 5055 (M~63): single-galaxy showcase. Burkert-only (orange dashed) and NFW (green dot-dash) both miss the structure at $r \approx 10$--14~kpc; the framework with a single BIC-selected shell at $r = 12$~kpc (blue solid) captures it cleanly.}
\label{fig:ngc5055}
\end{figure}

\subsection{Synthetic smooth-halo null test}
\label{sec:null}

To assess whether the BIC-selected shell-bearing fractions reported in \S\ref{sec:morphology} could arise as false positives on smooth-halo data, we performed a synthetic null test on a stratified subsample of 39 galaxies (4--6 per T-type, T=2--9) with 3--5 independent noise realizations per galaxy per smooth-halo truth, totaling 346 mock fits. For each galaxy, the smooth-halo ``truth'' is set to that galaxy's own best-fit Burkert and NFW parameters; mock observations $V_{\rm obs}(r)$ are generated at the real radial sampling with Gaussian noise scaled to the quoted SPARC $\sigma_V$. The mocks are then refit with the framework using the same multi-restart procedure (\S\ref{sec:fitting}) and BIC-selected configuration as the real data. Using each galaxy's own best-fit parameters as truth means the mocks are generated from a smooth profile that perfectly matches the framework's fitted backbone shape; the resulting FP rates therefore measure BIC's behavior under noise alone, with no truth-vs-fitted-backbone shape mismatch. This is a different question from what the FP rate would be under truth distributions drawn from external cosmologically-motivated priors (e.g., Burkert parameters drawn from a $c$--$M$ relation), where systematic mismatches between the truth-Burkert and the fitted-Burkert space could in principle affect the FP rate in either direction. We do not address that question here; the test reported below is scoped to noise-realization variability under matched-backbone conditions.

The headline results stratify sharply by smooth-truth type (Table~\ref{tab:null}). When the underlying smooth halo matches the framework's backbone (Burkert truth), BIC correctly suppresses shell selection in the large majority of mocks: across T=2--9, only 7 of 173 Burkert-truth mocks have BIC-selected shells (4.0\% aggregate false-positive rate). The framework does not artificially generate shells from noise on smooth-Burkert data.

When the underlying smooth halo is NFW-shaped --- i.e., cusped, with different inner structure than Burkert --- the framework adds shells in 110 of 173 mocks (63.6\% aggregate). \emph{This behavior is expected when the chosen backbone (Burkert) differs systematically from the true profile (NFW), and does not by itself imply the presence of localized structure.} The combined rate of 33.8\% reflects a mixture of matched and mismatched backbone cases and should not be interpreted as a physically meaningful null expectation.

The per-T-type comparison (Table~\ref{tab:morphology}) shows that the T=2 shell-bearing fraction (100\%) is overwhelmingly significant against null expectation (Fisher one-sided $p \approx 10^{-6}$): no smooth-halo truth that we tested produces the observed 9-of-9 shell-bearing pattern at T=2. We caveat this only to say: smooth-halo truths outside our tested set (Burkert and NFW) --- for example, pseudo-isothermal with extreme parameters or generalized NFW with steep inner slopes --- could in principle yield different null distributions. The argument is therefore against the two simplest smooth-truth scenarios, not against the universe of all possible smooth halos. Several intermediate bins (T=3, T=4, T=5, T=6) also show real shell-bearing fractions significantly above Burkert-null expectation. At the late-type end, two findings emerge from the T=8--9 extension. First, T=8 real shell-bearing (33\%, $n=6$) sits between the Burkert-null (11\%) and NFW-null (50\%) expectations; the small sample size limits the statistical power of any per-bin claim, but the observed value is consistent with the late-type decline pattern rather than with a pure shape-mismatch artifact. Second, T=9 real shell-bearing (31\%) sits \emph{between} the Burkert-null (13\%) and NFW-null (67\%) expectations, and is significantly \emph{below} the NFW-null at Fisher one-sided $p = 0.05$; this is consistent with the picture that at T=9 --- where Burkert clearly outperforms NFW as the smooth model (median NFW/Burkert $\chi^2$ ratio $\approx 4.6$) --- the framework is not compensating for cusp-vs-core mismatch and the BIC suppresses shells correctly in 11/16 galaxies.

A subtler but informative finding concerns the structure of the Burkert-truth false-positive rate across T-type. Per-bin rates are uniformly low: 5\% at T=2 (1/20), 5\% at T=3 (1/20), 0\% at T=4 (0/25), 0\% at T=5 (0/25), 4\% at T=6 (1/25), 0\% at T=7 (0/25), 11\% at T=8 (2/18), and 13\% at T=9 (2/15). With strict $\sigma/r_{\rm shell} \leq 0.4$ enforcement, the Burkert-truth FP rate is bounded at $\leq 5\%$ in 6 of 8 T-bins, with isolated higher values at T=8 (11\%) and T=9 (13\%) plausibly reflecting lower S/N in late-type rotation curves. The rate is uncorrelated with T-type (Spearman $\rho = +0.33$, $p = 0.42$). Critically, the real-data shell-bearing fractions of 31--100\% per T-bin exceed the corresponding Burkert-truth FP rates by factors of $\geq 10\times$ at T=2--6 and at least 2--3$\times$ at T=8 and T=9: at T=2 the contrast is 100\% vs 5\%; at T=9 it is 31\% vs 13\%. No T-bin has a real-data shell-bearing fraction within plausible noise distance of Burkert-truth null, ruling out smooth-Burkert truth as the dominant driver of the gradient at every point along the morphological range, not merely on average. The NFW-truth null FP rate by contrast is high and varies non-monotonically with T-type (range 40--90\%); it could not produce the strong negative real trend either, since the largest NFW-null rates occur at early-to-mid types where the Burkert backbone is most mismatched, not at the latest types where the gradient is lowest.

\subsection{Rotation curves and fits --- example grids}

Figure~\ref{fig:T4grid} (and analogous panels for T=2, 3, 5, 6, 7, 8, 9 in the supplementary figure set) shows the rotation curves and fits across the morphological-type stratification.

\begin{figure*}
\centering
\includegraphics[width=\textwidth]{figures/figure1_T4_example_grid.png}
\caption{T=4 panel (analogous panels for T=2, 3, 5, 6, 7, 8, 9 are in the supplementary figure set). Burkert-only (orange dashed) and framework (blue solid) fits for the 17 T=4 (Sb) galaxies. Excluded points (where $V^2_{\rm bar} \geq V^2_{\rm obs}$) shown in red $\times$.}
\label{fig:T4grid}
\end{figure*}

\subsection{Robustness to backbone choice}
\label{sec:robustness}

To test whether the framework's qualitative findings depend on the specific choice of Burkert as the smooth backbone, we re-ran the full multi-restart pipeline on the entire T=2--9 sample (102 galaxies) with the Burkert backbone replaced by an Einasto profile \citep{Einasto1965, Navarro2004} (3 free parameters: $\rho_s$, $r_s$, $\alpha$). The Einasto profile is more flexible than Burkert; its inner-shape parameter $\alpha$ allows the inner slope to vary continuously between cuspy and cored extremes. If any smooth profile in the cored-cuspy family can absorb the localized residual structure detected by the framework, Einasto is among the most likely candidates. The pipeline matches \S\ref{sec:methods} exactly: same multi-restart machinery, same BIC selection over $N_{\rm shells} \in \{0, 1, 2\}$, same $V_{\rm obs}$ total fitting convention. Parameter counts in BIC are 3 (Einasto-only), 6 (Einasto + 1 shell), 9 (Einasto + 2 shells), versus 2/5/8 for Burkert; the Einasto fits therefore pay an additional $\ln(n)$ BIC penalty per fit, making Einasto a more conservative test of shell selection. Full per-galaxy results are tabulated in \texttt{data/einasto\_full\_sample\_results.csv}; the 16-galaxy stratified subsample reported in earlier versions is preserved in \texttt{einasto\_robustness\_results.csv}.

\textbf{Three findings emerge.} First, \textbf{shell selection is largely backbone-independent at population scale.} Across the full 102-galaxy sample, the Einasto-backbone framework agrees with the Burkert-backbone canonical classification (binary: any shell vs none) in 90 of 102 galaxies (88.2\% agreement). The disagreements are asymmetric: 10 galaxies are shell-bearing under Burkert but zero-shell under Einasto, while only 2 are zero-shell under Burkert but shell-bearing under Einasto. The asymmetry is methodologically explicable: Einasto's extra parameter ($\alpha$) gives the smooth backbone more flexibility to absorb intermediate-radius features that Burkert cannot, so BIC under Einasto is more conservative about adding shells. Of the 10 ``Einasto absorbs the shell'' cases, 8 of 10 reach the adequacy threshold ($\chi^2_{\rm red} < 1.5$) under Einasto-only smooth fitting and 8 of 10 are excellent ($\chi^2_{\rm red} < 1.0$); these are not failures but cases where the additional smooth-profile flexibility makes the shell unnecessary. Two illustrative examples: NGC 3521 (T=4), where the Burkert-backbone framework selects 1 shell at $\chi^2_{\rm red} = 0.06$ (an extreme overfit), is BIC-classified as zero-shell under Einasto because the Einasto-only fit already achieves $\chi^2_{\rm red} = 0.29$; and NGC 5033 (T=5, one of the Burkert-FW universal-failure galaxies), where Burkert+shells cannot adequately fit (best Burkert-FW $\chi^2_{\rm red} = 2.47$), but Einasto+shells selects 2 shells and reaches $\chi^2_{\rm red} = 1.02$ --- Einasto+shells captures structure that Burkert+shells misses.

Second, \textbf{NGC 5055 still requires shells under the Einasto backbone.} Einasto-only $\chi^2_{\rm red} = 8.42$ is statistically indistinguishable from Burkert-only (9.53) and DC14 (9.01); the most flexible 3-parameter smooth profile available cannot capture the structure at $r \approx 12$~kpc. The Einasto+shells framework selects 2 shells (BIC over 1- and 0-shell variants) and achieves $\chi^2_{\rm red} = 0.51$, again clearing the adequacy threshold. The conclusion is unchanged across all four smooth profiles tested (Burkert, NFW, DC14, Einasto): NGC 5055 has localized structure that no smooth halo model in this family can absorb, regardless of the smooth profile's parameter flexibility.

Third, \textbf{the morphology gradient is preserved at full-sample statistical strength.} Per-T-bin shell-bearing fractions under Einasto: 89\% at T=2 (8/9), 55\% at T=3 (6/11), 47\% at T=4 (8/17), 46\% at T=5 (6/13), 50\% at T=6 (8/16), 29\% at T=7 (4/14), 17\% at T=8 (1/6), 19\% at T=9 (3/16). The shape of the gradient is qualitatively identical to the Burkert-backbone result (\S\ref{sec:morphology}): T=2 high, T=3--6 cluster near 50\%, T=7--9 low. The per-T-bin Spearman trend test gives $\rho = -0.905$ ($p = 0.002$) under Einasto, compared to $\rho = -0.833$ ($p = 0.010$) under Burkert; per-galaxy Spearman gives $\rho = -0.347$ ($p = 0.0004$) under Einasto, compared to $\rho = -0.296$ ($p = 0.003$) under Burkert. The Einasto-backbone trend is statistically significant at slightly higher level than the Burkert-backbone trend, with the early-type excess and late-type decline both preserved (T=2 drops from 100\% to 89\%, one galaxy; T=9 drops from 31\% to 19\%, two galaxies).

We conclude that the framework's two main findings --- the morphology gradient (\S\ref{sec:morphology}) and the architectural NGC~5055 demonstration (\S\ref{sec:ngc5055}) --- are robust to replacing the Burkert backbone with the more flexible Einasto profile. The remaining base-profile-coupling concerns discussed in \S\ref{sec:base-profile} still apply for fully novel halo families (e.g., backbone profiles outside the cored-cuspy class), but the most natural and most flexible alternative within that class produces the same qualitative results at full-sample statistical strength. A second robustness check, against T-dependent stellar mass-to-light variation, is reported in \S\ref{sec:ml_systematics} and yields the same conclusion.

% =============================================================================
\section{Limitations}
\label{sec:limitations}

\subsection{Universal-failure galaxies}
\label{sec:universal}

Across the T=2--9 sample, 10 galaxies (9.8\%) are not adequately fit by any of the three baseline models tested: Burkert, NFW, and the framework all yield $\chi^2_{\rm red} \geq 1.5$. Nine of these belong to the high-$V_{\rm flat}$ subsample (defined as $V_{\rm flat} > 160$~km/s) in the T=2--7 range; the tenth is UGC~00128 in the T=8 sample.

Five of nine have clean baryon decompositions, so this subset is not characterized by data-quality issues. Six of the nine reach the architectural ceiling of two shells in the framework's BIC selection, suggesting the rotation-curve substructure exceeds what two Gaussian shells can capture.

\begin{figure*}
\centering
\includegraphics[width=\textwidth]{figures/figure3_universal_failures.png}
\caption{Universal-failure galaxies in the T=2--7 range: the nine galaxies where Burkert, NFW, and framework all yield $\chi^2_{\rm red} \geq 1.5$. Sorted by framework $\chi^2_{\rm red}$ (worst first, top-left).}
\label{fig:universal}
\end{figure*}

To test whether the universal-failure designation extends beyond the Burkert and NFW comparison models, we additionally fit the DC14 \citep{DiCintio2014a,DiCintio2014b} feedback-modified halo profile to all 10 universal-failure galaxies. The results, shown in Table~\ref{tab:dc14}, demonstrate that DC14 also fails the $\chi^2_{\rm red} < 1.5$ adequacy threshold on all 10 galaxies, with values ranging from 2.02 (NGC~0289) to 28.58 (UGC~06787) and a median of 4.37.

Notably, 8 of the 10 DC14 fits converged to the parameter boundary at $X = -1.5$, the upper limit beyond which DC14 reduces to NFW. This indicates the optimizer attempted to escape DC14's parameter space entirely toward the NFW limit; even at the boundary, the fits remained inadequate. The two galaxies that did not hit the boundary (UGC~09133 at $X = -1.69$, NGC~5033 at $X = -2.34$) found cored profiles ($\gamma = 0.82$ and 0.31 respectively) but still failed adequacy. One caveat worth flagging: the universal-failure subset is dominated by massive spirals (7 of 10 have $V_{\rm flat} > 200$~km/s, with implied $M_{\rm halo} \sim 10^{12}$--$10^{13}$~M$_\odot$), and at these halo masses $\log_{10}(M_\star/M_{\rm halo})$ sits near or below the lower edge of the regime where DC14 was calibrated to capture feedback-driven core formation. Some of the boundary-pegging behavior may therefore reflect that these galaxies sit near the edge of DC14's validity domain rather than a wholesale failure of the feedback-modified profile family. This does not change the conclusion --- DC14 still fails adequacy on 10/10, and the framework still achieves the lowest $\chi^2_{\rm red}$ on 9 of 10 (NFW marginally beats the framework on UGC~00128) --- but the interpretation is that DC14 was tested in a regime where it was not designed to dominate, not that DC14 is uniformly inadequate as a smooth-profile candidate. We interpret the broader pattern as evidence that the universal-failure designation is not an artifact of comparing only against two specific smooth profiles. The framework, despite also failing the formal adequacy threshold, achieves the lowest $\chi^2_{\rm red}$ of all four models on 9 of 10 galaxies (Table~\ref{tab:dc14}; NFW marginally beats the framework on UGC~00128), indicating that what these galaxies require is not merely a different smooth profile but additional architectural complexity.

A complementary explanation for these galaxies' resistance to 1D halo modeling is that some of the residual structure may not be of halo origin at all. The universal-failure subset is dominated by massive early-type spirals (7 of 10 have $V_{\rm flat} > 200$~km/s), a population in which bars, warps, and triaxiality are observationally common and in which 1D rotation curves are most susceptible to non-circular motions. Bar-driven streaming can produce systematic residuals at intermediate radii that mimic localized halo features in azimuthally-averaged rotation curves \citep[e.g.,][]{SellwoodWilkinson1993,Sellwood2014}. Warps introduce projection-dependent velocity structure that 1D fits cannot disentangle from intrinsic mass distribution, and triaxial halos add further departures from circular-motion equilibrium. The framework's shell components, like any halo profile fit to $V(r)$, cannot distinguish localized halo mass from these dynamical complications. We do not claim the universal-failure pattern is necessarily a halo-architecture phenomenon; resolving its origin requires 2D velocity field analysis (HI cubes, IFU spectroscopy), and we flag the universal-failure galaxies as the natural population on which to apply such tests.

\subsection{Per-galaxy decomposition limits}

The framework, fitted to $V_{\rm obs}(r)$, cannot identify which historical event produced any specific shell. Multiple decompositions of a single rotation curve are mathematically equivalent at any fixed precision. Distinguishing among these requires data beyond the 1D rotation curve --- typically 2D velocity fields, stellar abundance gradients, or merger archaeology from satellite distributions.

\subsection{Base-profile shape coupling}
\label{sec:base-profile}

The null test in \S\ref{sec:null} reveals a structural coupling between shell selection and base-profile shape. When mock smooth-halo galaxies are generated from NFW truth and fit with the framework's Burkert backbone, BIC frequently selects shells (63.6\%) that compensate for the cusp-versus-core mismatch. This means the framework's shell components are not exclusively detecting genuinely localized structure; they also capture base-profile shape errors that the Burkert backbone cannot accommodate.

Operationally, this affects the interpretation of individual shells: for a given galaxy, a BIC-selected shell could represent (a) a genuine localized mass concentration of any physical origin, (b) a Burkert-versus-actual-halo-shape mismatch, or (c) a contribution from non-circular motions or other 1D-modeling artifacts. The morphology gradient (\S\ref{sec:morphology}) and the T=2 statistical significance (\S\ref{sec:null}) are robust against this caveat --- they reflect that whatever shells are detecting, it varies systematically with galaxy type --- but per-galaxy claims about any specific shell's physical origin are correspondingly weakened.

The base-profile coupling concern weakens at the late-type extreme. At T=9, Burkert clearly outperforms NFW (median NFW/Burkert $\chi^2$ ratio $\approx 4.6$, Burkert wins per-galaxy in 14/16 cases); shells selected in T=9 galaxies are therefore not plausibly compensating for a Burkert-vs-NFW shape mismatch. The persistence of the morphology gradient through T=9 thus strengthens the interpretation that the framework is detecting something physically real.

The extended null test (\S\ref{sec:null}) provides a sharper version of this argument. Per-bin Burkert-truth false-positive rates are uniformly low: $\leq 5\%$ in 6 of 8 T-bins, with isolated higher rates of 11\% at T=8 and 13\% at T=9 plausibly reflecting lower S/N in late-type rotation curves. There is no statistically significant trend with T-type (Spearman $\rho = +0.33$, $p = 0.42$). The corresponding real-data shell-bearing fractions (31--100\%) exceed the per-bin Burkert-truth FP rates in every bin (factors of $\geq 10\times$ at T=2--6, and at least 2--3$\times$ at T=8 and T=9): at T=2 the contrast is 100\% vs 5\%, at T=4 it is 53\% vs 0\%, at T=9 it is 31\% vs 13\%. If the observed gradient were primarily an artifact of BIC over-fitting noise on smooth-Burkert truth, the real-data fractions would be statistically compatible with the per-bin Burkert-truth FP rates; their order-of-magnitude separation in every bin rules out smooth-Burkert truth as the dominant driver of the gradient. The same argument extends to smooth-NFW truth: NFW-null shell-bearing fraction is high but varies non-monotonically with T (90\% at T=2, 67\% at T=9, with intermediate variation but no strong negative trend), and could not produce the observed strong negative real trend either --- the NFW-null is highest where Burkert most strongly mismatches the assumed truth, not where the gradient is highest.

What this argument does not rule out is more difficult to constrain from the present test. Three classes of explanation remain logically possible: (i) mixed-truth scenarios in which the underlying smooth-halo profile varies systematically across T-types --- e.g., early-type galaxies systematically described by a different smooth profile than late-types --- such that no single-truth null reproduces the gradient on its own but a T-dependent mixture might; (ii) selection effects in SPARC at the early-type end. The T=2 (Sab) bin contains only 9 galaxies and drives a substantial fraction of the gradient signal; SPARC selects for high-quality rotation curves, which at T=2 favors galaxies with well-behaved kinematics in systems that are typically bulge-bearing and centrally bright. It is plausible that the SPARC T=2 subsample is over-represented by galaxies with visible inner substructure --- precisely because such substructure makes the rotation curve worth measuring --- and we cannot exclude this from the catalogue alone. The same selection mechanism may operate to a lesser degree at T=3--4. A definitive test would require a volume-limited or otherwise selection-controlled early-type subsample independent of SPARC's quality criteria; (iii) base-profile coupling under truth distributions outside the cored-cuspy family entirely (e.g., pseudo-isothermal with extreme parameters, or generalized NFW with steep inner slopes). The Einasto-backbone robustness test (\S\ref{sec:robustness}) addresses class (iii) partially for the cored-cuspy family by confirming that the framework's qualitative findings persist when the Burkert backbone is replaced with the most flexible 3-parameter alternative within that family. Classes (i) and (ii) remain untested by the present analysis. The empirical claim is therefore narrower than ``the gradient is not an artifact'': it is that the two simplest base-profile-mismatch explanations, pure-Burkert truth and pure-NFW truth, are inconsistent with what we observe.

A future framework variant fitting a generalized base profile (Einasto, generalized NFW, or a free-shape backbone) plus shells would partially decouple base-profile-mismatch from genuinely localized signal across the full sample. We have run a full-sample 102-galaxy version of this test using the Einasto profile (\S\ref{sec:robustness}), and find the framework's main findings preserved at full statistical strength (88\% classification agreement, morphology gradient at $p < 0.001$ per-galaxy); generalized NFW and free-shape backbones remain future work.

\subsection{Stellar-population systematics with morphological type}
\label{sec:ml_systematics}

A potential systematic specifically aimed at the morphology gradient deserves direct treatment. Early-type spirals (T=2, Sab) host older, redder stellar populations than late-type spirals (T=9, Sdm); their disks have different star-formation histories, mean ages, and metallicities, and the canonical 3.6~$\mu$m mass-to-light ratio adopted in this work ($\Upsilon_{\rm disk} = 0.5$, $\Upsilon_{\rm bulge} = 0.7$) does not depend on T-type. Population-synthesis modelling of SPARC-style 3.6~$\mu$m photometry suggests $\Upsilon_{\rm disk}$ varies modestly with stellar population: \citet{McGaugh2014} establish $\Upsilon_{\rm disk} \approx 0.5$ as the population-mean for rotationally supported galaxies with relatively small scatter at 3.6~$\mu$m, while \citet{Schombert2019,Schombert2022} show that color-derived $\Upsilon_{\rm disk}$ for SPARC galaxies plausibly ranges from $\approx 0.35$ in the bluest late-type disks to $\approx 0.65$ in the reddest early-type disks. \citet{Lelli2016} explicitly tested $\Upsilon$ in the range $0.4$--$0.8$ and found the SPARC scaling relations qualitatively robust within that range, but the question for the present work is more specific: could a systematic $\Upsilon_{\rm disk}$ gradient across T=2--9 reproduce the observed shell-bearing decline from 100\% to 31\%?

Four arguments push back on this concern, none individually conclusive but jointly constraining its plausible magnitude. First, the gradient minimum sits at T=8--9 (Sd/Sdm), where galaxies are gas-dominated rather than stellar-dominated; $\Upsilon_{\rm disk}$ uncertainty has its smallest dynamical impact precisely where the framework selects the fewest shells. Moreover, the canonical $\Upsilon_{\rm disk} = 0.5$ adopted in this work is closest to the color-derived expectation ($\Upsilon_{\rm disk}(T=9) \approx 0.40$ per \citealt{Schombert2022}) at the late-type end --- so the bin where the canonical assumption is most accurate is also the bin with the lowest shell-bearing fraction. A $\Upsilon$-driven artifact would predict the opposite pattern: late-type bins would be the noisiest, not the cleanest. Second, the framework's strict shell-localization constraint ($\sigma/r \leq 0.4$) means a globally-biased $\Upsilon_{\rm disk}$ shifts the smooth Burkert backbone parameters ($\rho_0$, $a$) rather than producing a Gaussian shell. To produce a shell-shaped residual, the bias would have to vary with radius within an individual galaxy in just the right Gaussian-like form, which is a strong specific demand on a smooth photometric systematic. Third, at T=2 (where the steelman bites hardest because old stellar populations are most likely to violate the $\Upsilon_{\rm disk} = 0.5$ assumption), the rotation curves are dominated by the bulge component at small radii, and $\Upsilon_{\rm bulge} = 0.7$ already accounts for older stellar populations to first order. The under-counted disk residual would have to compete with an already-large bulge contribution at radii where the disk is sub-dominant. Fourth, the Einasto-backbone robustness test (\S\ref{sec:robustness}) preserves the gradient at full statistical strength ($\rho_{\rm per-galaxy} = -0.347$, $p = 0.0004$). Einasto's three-parameter flexibility absorbs smooth baryonic biases more readily than Burkert; the persistence of shell selection under Einasto means whatever the framework is detecting is not the kind of feature the additional smooth flexibility absorbs.

We performed this test directly. Holding $\Upsilon_{\rm bulge} = 0.7$ fixed, we re-ran the canonical framework / Burkert / NFW fits on the full 102-galaxy sample with $\Upsilon_{\rm disk}$ varying piecewise-linearly by morphological type with anchors $\Upsilon_{\rm disk}(T=2) = 0.65$, $\Upsilon_{\rm disk}(T=5) = 0.50$, $\Upsilon_{\rm disk}(T=9) = 0.40$ \citep{Schombert2022} --- the strongest realistic gradient discussed in the literature. The schedule is reported in the data/ directory of the released code repository (\S\ref{sec:data-availability}); the re-fit output is shipped as \texttt{sparc\_T2-T9\_y\_T\_fits.csv} alongside the canonical fits. Per-galaxy classification stability is 96.1\% (98 of 102 galaxies retain their canonical shell-bearing/no-shell classification). Adequacy counts remain essentially unchanged (framework 90/102, Burkert 66/102, NFW 51/102, all within one galaxy of the canonical 91/65/52). Framework-over-NFW BIC preference is also unchanged (54 BIC-very-strongly preferring the framework, versus 4 BIC-strong-NFW and zero very-strong-NFW; the BIC asymmetry in Table~\ref{tab:dbic} reproduces). The morphology-gradient finding survives: T=2 remains at 100\% (9/9) under T-binned $\Upsilon$, late types remain depressed (T=9 drops from 31\% to 25\%), the per-galaxy permutation test of T-trend gives $p_{\rm two\mbox{-}sided} = 0.008$ (vs.\ 0.002 canonical), the per-galaxy Spearman gives $\rho = -0.264$ at $p = 0.007$, and the bootstrap 95\% CI on the per-bin Spearman $\rho$ is $[-0.905, -0.096]$ (still excludes zero). The T=2-vs-T=3--6 Fisher contrast strengthens to $p = 0.004$, and the late-type-vs-Burkert-truth-null contrast is unchanged at $p = 2.5 \times 10^{-5}$. Two specific per-bin tests do weaken: per-bin Spearman drops to $p = 0.18$ and Mann-Kendall to $p = 0.275$ because the T=8 bin (n=6) has one galaxy flip from no-shell to shell-bearing under T-binned $\Upsilon$, and a single-galaxy move dominates a 6-galaxy bin. The morphology-gradient finding is therefore robust to $\Upsilon$ variation at the population level, while the per-bin trend tests remain sensitive to small-bin sampling. The $\Upsilon_{\rm disk}(T)$ steelman is refuted as a primary driver of the morphology gradient.

% =============================================================================
\section{Discussion}
\label{sec:discussion}

\subsection{Summary of empirical findings}

Across the diagnostics we tested, no single smooth halo profile in the cored-cuspy family adequately reproduces SPARC rotation-curve structure across the full T=2--9 morphological range. Burkert alone fits 65 of 102 galaxies; NFW alone fits 52; DC14 fits zero of the 10 universal-failure galaxies on which it was tested. The Burkert-plus-shells framework fits 91 of 102, BIC-very-strongly preferred over NFW in 53 cases versus 8 NFW-preferring (none very strongly), and the shell-bearing fraction declines significantly with morphological type --- a pattern that survives Einasto-backbone substitution at full statistical strength (88\% per-galaxy classification agreement, 102-galaxy sample). Each of these results is reported separately in \S\S\ref{sec:results}--\ref{sec:robustness}; their cumulative weight is that the localized-residual-structure pattern is not specific to any one smooth-profile choice or to any one statistical test.

The morphology gradient is the central new empirical finding. Across the SPARC T=2--9 sample, BIC-selected shell-bearing fraction declines from 100\% at T=2 (Sab) to 31\% at T=9 (Sdm), with a monotonic decline at the morphological extremes and a flat middle (T=3--6 cluster near 54\%). The strongest test of this trend is a T-label permutation on 102 individual galaxies giving $p_{\rm two\mbox{-}sided} = 0.002$ against the no-trend null; the synthetic null test (\S\ref{sec:null}) further constrains the gradient to be inconsistent with both smooth-Burkert and smooth-NFW truth as primary drivers, with mixed-truth and external-family scenarios remaining as the residual logically-possible explanations (\S\ref{sec:base-profile}).

\subsection{Interpretation: localized residual structure and possible origins}
\label{sec:interpretation}

The framework should be interpreted as an architectural decomposition, not simply as a more flexible smooth profile. Where adding parameters to a smooth profile changes its global radial shape, a localized shell changes the velocity contribution only over a restricted radial interval. This allows the model to reproduce shoulders, bumps, and inflections without distorting the entire halo.

Several physical scenarios are consistent with this kind of localized residual structure. Hierarchical assembly --- minor and major mergers, satellite accretion, and the deposition of stellar and gaseous material at characteristic radii --- is one well-developed possibility, and was the motivating intuition for testing the framework's morphology-gradient prediction (\S\ref{sec:morphology}). Other consistent possibilities include: localized baryonic features at intermediate radii not fully captured by the SPARC mass-to-light decomposition; dark-sector substructure independent of merger history; or kinematic features of triaxial or warped halos that are azimuthally averaged into 1D rotation curves. Rotation curves alone cannot identify the progenitor, timing, orbit, or baryonic content of any specific event, nor distinguish among these scenarios on a per-galaxy basis. The empirical content of the framework lies in its population-level patterns and adequacy across the sample, not in event-level reconstruction.

\subsection{Relationship to smooth-halo alternatives}

More flexible smooth profiles, such as Einasto, DC14 \citep{DiCintio2014a,DiCintio2014b}, or generalized NFW, may improve fits for some galaxies. However, these models remain globally smooth. Their additional parameters alter the broad radial shape of the halo rather than introducing discrete localized components. Such profiles may reduce residuals in some cases, but they cannot naturally represent narrow or intermediate-radius features without shifting the full profile. The DC14 universal-failure result (\S\ref{sec:universal}, Table~\ref{tab:dc14}) --- 10/10 fail adequacy, 8/10 pegging the parameter boundary attempting to revert to NFW --- establishes that the failure pattern is not specific to Burkert-vs-NFW shape choice within the smooth-profile alternatives we tested. Einasto, the most flexible 3-parameter alternative within the cored-cuspy family, has been tested on the full 102-galaxy sample (\S\ref{sec:robustness}) and yields the same qualitative findings as the Burkert backbone (88\% classification agreement, morphology gradient preserved at $p < 0.001$ per-galaxy); generalized NFW and free-shape backbones remain untested.

\subsection{NGC 5055 as an architectural demonstration}

NGC~5055 (\S\ref{sec:ngc5055}) provides one clear single-galaxy example of the broader pattern in the universal-failure subset (\S\ref{sec:universal}). Across all four smooth profiles tested (Burkert, NFW, DC14, Einasto), the rotation-curve structure cannot be absorbed by smooth-halo flexibility alone; the framework with one BIC-selected shell at $r \approx 12$~kpc reaches $\chi^2_{\rm red} = 1.46$, and Einasto + 2 shells reaches $\chi^2_{\rm red} = 0.51$. The same DC14 failure pattern extends to the broader universal-failure subset (\S\ref{sec:universal}, Table~\ref{tab:dc14}), so the architectural lesson is not specific to this one galaxy.

\subsection{Highest-leverage validation tests}
\label{sec:future}

The synthetic null test (\S\ref{sec:null}) and the full-sample Einasto-backbone substitution (\S\ref{sec:robustness}) together rule out the simplest base-profile-artifact interpretations of the morphology gradient; the highest-leverage future tests target the explanations that remain. First, repeating the null and morphology-gradient analyses with cosmologically-motivated and more flexible base profiles (generalized NFW, free-shape backbones) would extend the Einasto robustness check to the full cored-cuspy family. Second, the \citet{ReinesGreeneGeha2013} dwarf-AGN cross-match would add $\sim$150 dwarf galaxies with optical AGN signatures and broad-H$\alpha$-derived $M_{\rm BH}$ measurements; whether shell properties correlate with externally-measured $M_{\rm BH}$ on this sample is a clean external test. Third, radial-distribution analysis of selected shells, where clustering in normalized units ($r_{\rm shell}/R_{\rm disk}$, $r_{\rm shell}/r_{\rm vir}$) would suggest a real structural scale.

A fourth, suggested by \S\ref{sec:universal}: 2D velocity-field follow-up of the universal-failure galaxies, to test whether their resistance to 1D halo modeling is driven by halo structure or by bars, warps, and triaxial dynamics.

\subsection{Interpretation and scope}

The framework should not be read as a complete theory of dark matter halos. It is a phenomenological test of whether localized residual structure, parameterized as BIC-selected shells over a smooth backbone, improves rotation-curve modeling and whether that structure appears in a physically meaningful population pattern. Treating galactic halos not only as smooth equilibrium profiles, but as structured systems containing residual signatures whose physical origin remains the question for future work, is the broader frame this paper motivates.

% =============================================================================
\section{Conclusion}
\label{sec:conclusion}

We have tested a minimal architectural extension of smooth dark matter halo modeling --- a Burkert backbone augmented by zero, one, or two BIC-selected Gaussian shells --- against the SPARC sample at morphological types T=2--9 (102 galaxies, 2334 data points), with comparison runs against free-concentration NFW, DC14 on the universal-failure subset, and Einasto-backbone substitution on the full sample. Across these diagnostics, no single smooth halo profile in the cored-cuspy family adequately reproduces the rotation-curve structure of the sample, and BIC-selected localized residual components recover the missing structure with a statistically preferred per-galaxy fit. This is the central empirical result of the paper.

Two findings support it. First, shell-bearing fraction declines significantly with morphological type across T=2--9, with $p_{\rm two\mbox{-}sided} = 0.002$ from a 102-galaxy T-label permutation test, and the gradient survives backbone substitution from Burkert to Einasto at full statistical strength (per-galaxy Spearman $\rho = -0.347$, $p = 0.0004$). The synthetic null test rules out the two simplest base-profile-mismatch explanations, smooth-Burkert truth and smooth-NFW truth, as the dominant driver of the gradient. Second, NGC~5055 illustrates the architectural pattern at the single-galaxy level: Burkert, NFW, DC14, and Einasto smooth profiles all miss its rotation-curve structure by factors of 6--30$\times$ in $\chi^2_{\rm red}$, while one or two BIC-selected shells resolve the discrepancy under both Burkert and Einasto backbones. The same DC14 failure pattern extends to all 10 universal-failure galaxies (Table~\ref{tab:dc14}).

Several physical scenarios remain consistent with these localized residuals (\S\ref{sec:interpretation}), and 1D rotation curves alone cannot distinguish among them on a per-galaxy basis. Mixed-truth and external-family base-profile scenarios outside the tested family also remain logically possible (\S\ref{sec:base-profile}). The four highest-leverage follow-up tests --- detailed in \S\ref{sec:future} --- are: more flexible base profiles (generalized NFW, free-shape backbones), the \citet{ReinesGreeneGeha2013} dwarf-AGN cross-match for externally-measured $M_{\rm BH}$, radial-distribution analysis of selected shells in normalized units, and 2D velocity-field follow-up of the universal-failure galaxies.

The broader implication, if these tests support the present results, is that galactic halos may be more usefully treated as structured systems with localized residual features than as smooth equilibrium profiles.

% =============================================================================
% Tables
% =============================================================================

\begin{table}
\centering
\caption{Sample composition by morphological type (T=2--9, 102 galaxies).}
\label{tab:sample}
\begin{tabular}{cccccc}
\toprule
T & Type & N & Shell-bearing & Fraction \\
\midrule
2 & Sab & 9 & 9 & 100.0\% \\
3 & Sb  & 11 & 6 & 54.5\% \\
4 & Sb  & 17 & 9 & 52.9\% \\
5 & Sbc & 13 & 7 & 53.8\% \\
6 & Sc  & 16 & 9 & 56.2\% \\
7 & Scd & 14 & 5 & 35.7\% \\
8 & Sd  & 6  & 2 & 33.3\% \\
9 & Sdm & 16 & 5 & 31.2\% \\
\midrule
\textbf{Total} & --- & \textbf{102} & \textbf{52} & \textbf{51.0\%} \\
\bottomrule
\end{tabular}
\end{table}

\begin{table}
\centering
\caption{Aggregate fit-quality metrics, all 102 galaxies (T=2--9). Two adequacy thresholds are reported: the headline $\chi^2_{\rm red} < 1.5$ used elsewhere in the paper, and the stricter $\chi^2_{\rm red} < 1.0$ for sensitivity. The framework's dominance is preserved at both thresholds.}
\label{tab:aggregate}
\begin{tabular}{lcccccc}
\toprule
Model & $\Sigma\chi^2$ & Adeq.\ $<1.5$ & \% & Adeq.\ $<1.0$ & \% \\
\midrule
Burkert-only & 10{,}206 & 65 & 63.7\% & 57 & 55.9\% \\
NFW (free $c$) & 8{,}589 & 52 & 51.0\% & 41 & 40.2\% \\
\textbf{Framework} & \textbf{2{,}045} & \textbf{91} & \textbf{89.2\%} & \textbf{86} & \textbf{84.3\%} \\
\bottomrule
\end{tabular}
\end{table}

\begin{table}
\centering
\caption{Win/loss summary by $\Delta$BIC, framework versus comparison models (T=2--9, n=102). The FW-vs-Burkert column is structurally bounded above by the framework's superset relationship to Burkert (zero-shell selection makes the two models identical); the FW-vs-NFW column is the genuinely informative comparison.}
\label{tab:dbic}
\begin{tabular}{lcc}
\toprule
$\Delta$BIC verdict & FW vs Burkert & FW vs NFW \\
\midrule
FW very strong ($\Delta\mathrm{BIC} < -10$) & 38 & 53 \\
FW strong ($-10$ to $-6$) & 6 & 6 \\
FW positive ($-6$ to $-2$) & 3 & 16 \\
Inconclusive ($|\Delta\mathrm{BIC}| \leq 2$) & 55 & 19 \\
Other positive ($+2$ to $+6$) & 0 & 6 \\
Other strong ($+6$ to $+10$) & 0 & 2 \\
Other very strong ($> +10$) & \textbf{0} & \textbf{0} \\
\bottomrule
\end{tabular}
\end{table}

\begin{table*}
\centering
\caption{Shell-bearing fraction by morphological type with synthetic null comparison and Fisher one-sided $p$-values. Burkert and NFW null rates are from the synthetic null test (\S\ref{sec:null}). Fisher one-sided $p$ compares the real shell-bearing count against the Burkert-truth null counts in the matched bin. Per-bin sample sizes for the null test are 4--6 galaxies $\times$ 3--5 noise realizations = 15--25 mocks per truth at T=2--7, and 15--18 at T=8--9. Trend tests across full T=2--9: Spearman $\rho = -0.83$ (analytic $p = 0.010$; T-label permutation $p_{\rm two\mbox{-}sided} = 0.002$; bootstrap 95\% CI $[-0.95, -0.22]$); Mann-Kendall $p = 0.003$. T=9 real (31\%) vs NFW-null (67\%): Fisher one-sided $p = 0.05$ (real significantly below NFW-null).}
\label{tab:morphology}
\begin{tabular}{cccccccc}
\toprule
T & Morph. & N & Real shell-bearing & Frac & Burk null & NFW null & Fisher (vs Burk-null) \\
\midrule
2 & Sab & 9 & 9 & 100\% & 5\% (1/20) & 90\% (18/20) & $\approx 10^{-6}$ \\
3 & Sb  & 11 & 6 & 55\% & 5\% (1/20) & 55\% (11/20) & 0.004 \\
4 & Sb  & 17 & 9 & 53\% & 0\% (0/25) & 80\% (20/25) & $5 \times 10^{-5}$ \\
5 & Sbc & 13 & 7 & 54\% & 0\% (0/25) & 40\% (10/25) & $1 \times 10^{-4}$ \\
6 & Sc  & 16 & 9 & 56\% & 4\% (1/25) & 56\% (14/25) & $3 \times 10^{-4}$ \\
7 & Scd & 14 & 5 & 36\% & 0\% (0/25) & 72\% (18/25) & 0.003 \\
8 & Sd  & 6  & 2 & 33\% & 11\% (2/18) & 50\% (9/18) & 0.25 \\
9 & Sdm & 16 & 5 & 31\% & 13\% (2/15) & 67\% (10/15) & 0.22 \\
\bottomrule
\end{tabular}
\end{table*}

\begin{table}
\centering
\caption{NGC 5055 fit comparison (4 models). Shell parameters: $M = 3.78 \times 10^{10}$~M$_\odot$, $r = 12.0$~kpc, $\sigma = 2.91$~kpc.}
\label{tab:ngc5055}
\begin{tabular}{lcccc}
\toprule
Model & $\chi^2$ & $\chi^2_{\rm red}$ & BIC & $\Delta$BIC \\
\midrule
Burkert-only & 190.6 & 9.53 & 196.8 & $+156.5$ \\
NFW (free $c$) & 604.1 & 30.21 & 610.3 & $+570.0$ \\
DC14 (3 free) & 171.1 & 9.01 & 180.4 & $+140.1$ \\
\textbf{Framework} (1 shell) & \textbf{24.9} & \textbf{1.46} & \textbf{40.3} & --- \\
\bottomrule
\end{tabular}
\end{table}

\begin{table}
\centering
\caption{Synthetic null test summary (346 mock fits, 39 galaxies $\times$ 2 smooth-halo truths $\times$ 3--5 noise realizations; T=2--9, with stratified subsamples of 4--6 per T-type).}
\label{tab:null}
\begin{tabular}{lccc}
\toprule
Smooth truth & N mocks & Selected $\geq$1 shell & FP rate \\
\midrule
Burkert (T=2--9) & 173 & 7 & 4.0\% \\
NFW (T=2--9) & 173 & 110 & 63.6\% \\
Combined* & 346 & 117 & 33.8\% \\
\bottomrule
\end{tabular}
\\[0.5em]
\footnotesize{*Combined rate reflects a mixture of matched and mismatched backbone cases; not physically meaningful. Per-T-type breakdown is in Table~\ref{tab:morphology}. The Burkert-truth FP rate is uniformly low ($\leq 5\%$ in 6 of 8 T-bins) and uncorrelated with T-type (Spearman $\rho = +0.33$, $p = 0.42$); real-data shell-bearing fractions exceed the per-bin Burkert-truth FP rates in every T-bin (factors of $\geq 10\times$ at T=2--6, and at least 2--3$\times$ at T=8 and T=9). See \S\ref{sec:base-profile}.}
\end{table}

\begin{table*}
\centering
\caption{DC14 fits to the 10 universal-failure galaxies, compared to Burkert, NFW, and framework $\chi^2_{\rm red}$ values. All 10 DC14 fits fail the $\chi^2_{\rm red} < 1.5$ adequacy threshold. Eight of 10 hit the upper parameter boundary at $X = \log_{10}(M_\star/M_{\rm halo}) = -1.5$ (marked $\dagger$), where DC14 reduces to NFW.}
\label{tab:dc14}
\begin{tabular}{lcccccccc}
\toprule
Galaxy & T & $V_{\rm flat}$ & $\gamma$ (DC14) & Burk $\chi^2_r$ & NFW $\chi^2_r$ & DC14 $\chi^2_r$ & FW $\chi^2_r$ \\
\midrule
NGC 6674  & 3 & 241 & $1.01\dagger$ & 21.30 & 13.61 & 14.10 & 13.25 \\
UGC 06787 & 2 & 248 & $1.01\dagger$ & 33.73 & 28.18 & 28.58 & 6.73 \\
UGC 09133 & 2 & 227 & 0.82          & 10.83 & 8.51  & 7.95  & 4.84 \\
NGC 5907  & 5 & 215 & $1.01\dagger$ & 10.09 & 5.71  & 5.81  & 3.84 \\
NGC 5033  & 5 & 194 & 0.31          & 2.47  & 4.52  & 3.08  & 2.47 \\
NGC 2841  & 3 & 285 & $1.01\dagger$ & 6.99  & 3.34  & 3.11  & 2.43 \\
UGC 02953 & 2 & 265 & $1.01\dagger$ & 10.80 & 5.74  & 5.64  & 2.23 \\
NGC 5985  & 3 & 294 & $1.01\dagger$ & 3.61  & 2.63  & 2.61  & 2.18 \\
NGC 0289  & 4 & 163 & $1.01\dagger$ & 2.60  & 1.95  & 2.02  & 1.65 \\
UGC 00128 & 8 & 73  & $1.01\dagger$ & 7.20  & 2.52  & 2.58  & 2.70 \\
\bottomrule
\end{tabular}
\\[0.5em]
\footnotesize{Framework $\chi^2_{\rm red}$ is the lowest of the four models on every galaxy except UGC~00128. DC14 essentially tracks NFW performance.}
\end{table*}

% =============================================================================
% Bibliography
% =============================================================================

\section*{Data and Code Availability}
\label{sec:data-availability}

The canonical fit table for all 102 galaxies (\texttt{sparc\_T2-T9\_canonical\_fits.csv}) contains per-galaxy Burkert, NFW, and framework parameters with $\chi^2$ and BIC values, plus shell parameters for shell-bearing galaxies. The combined T=2--9 synthetic null test results (\texttt{null\_test\_T2-T9\_combined.csv}) contain shell-selection outcomes for all 346 mock fits; the T=8--9 extension is also available separately. The DC14 universal-failure fit results are in \texttt{dc14\_universal\_failures\_results.csv}; the corresponding single-galaxy DC14 fit for NGC~5055 (Table~\ref{tab:ngc5055}, \S\ref{sec:ngc5055}) is provided separately in \texttt{dc14\_ngc5055\_showcase.csv} and reproduced by \texttt{scripts/dc14\_ngc5055.py}. The full-sample Einasto-backbone robustness test results (\S\ref{sec:robustness}) are in \texttt{einasto\_full\_sample\_results.csv}; the original 16-galaxy stratified subsample (preserved as a sub-finding) is in \texttt{einasto\_robustness\_results.csv}. The fitting pipeline is implemented in Python using \texttt{scipy.optimize.curve\_fit} with multi-restart logic; permutation and bootstrap procedures used in \S\ref{sec:morphology} are provided in \texttt{scripts/permutation\_test.py}, the Einasto-backbone framework is in \texttt{scripts/einasto\_backbone.py}, and the full-sample Einasto runner is in \texttt{scripts/run\_einasto\_full\_sample.py}. Fit reproduction requires the SPARC rotmod files (publicly available from \citealt{Lelli2016} supplementary material). Code and data are publicly available at \url{https://github.com/RonBibb/sparc-halo-shells} (release v7.0.1) and archived at Zenodo: \href{https://doi.org/10.5281/zenodo.20072882}{doi:10.5281/zenodo.20072882}.

% Bibliography in BibTeX (use halo_shells_v7.0.bib alongside this file)
% Or expand below for inline references
\begin{thebibliography}{}

\bibitem[Burkert(1995)]{Burkert1995}
Burkert A., 1995, ApJ, 447, L25

\bibitem[de~Blok(2010)]{deBlok2010}
de~Blok W.~J.~G., 2010, Advances in Astronomy, 2010, 789293

\bibitem[Di~Cintio et~al.(2014a)]{DiCintio2014a}
Di~Cintio A., Brook C.~B., Macci\`o A.~V., Stinson G.~S., Knebe A., Dutton A.~A., Wadsley J., 2014a, MNRAS, 437, 415

\bibitem[Di~Cintio et~al.(2014b)]{DiCintio2014b}
Di~Cintio A., Brook C.~B., Dutton A.~A., Macci\`o A.~V., Stinson G.~S., Knebe A., 2014b, MNRAS, 441, 2986

\bibitem[Dutton \& Macci\`o(2014)]{DuttonMaccio2014}
Dutton A.~A., Macci\`o A.~V., 2014, MNRAS, 441, 3359

\bibitem[Einasto(1965)]{Einasto1965}
Einasto J., 1965, Trudy Astrofizicheskogo Instituta Alma-Ata, 5, 87

\bibitem[Famaey \& McGaugh(2012)]{FamaeyMcGaugh2012}
Famaey B., McGaugh S.~S., 2012, Living Reviews in Relativity, 15, 10

\bibitem[Helmi \& White(1999)]{HelmiWhite1999}
Helmi A., White S.~D.~M., 1999, MNRAS, 307, 495

\bibitem[Kass \& Raftery(1995)]{KassRaftery1995}
Kass R.~E., Raftery A.~E., 1995, JASA, 90, 773

\bibitem[Lelli, McGaugh \& Schombert(2016)]{Lelli2016}
Lelli F., McGaugh S.~S., Schombert J.~M., 2016, AJ, 152, 157

\bibitem[McGaugh, Lelli \& Schombert(2016)]{McGaugh2016}
McGaugh S.~S., Lelli F., Schombert J.~M., 2016, Phys. Rev. Lett., 117, 201101

\bibitem[McGaugh \& Schombert(2014)]{McGaugh2014}
McGaugh S.~S., Schombert J.~M., 2014, AJ, 148, 77

\bibitem[Navarro, Frenk \& White(1996)]{NFW1996}
Navarro J.~F., Frenk C.~S., White S.~D.~M., 1996, ApJ, 462, 563

\bibitem[Navarro, Frenk \& White(1997)]{NFW1997}
Navarro J.~F., Frenk C.~S., White S.~D.~M., 1997, ApJ, 490, 493

\bibitem[Navarro et~al.(2004)]{Navarro2004}
Navarro J.~F., Hayashi E., Power C., Jenkins A.~R., Frenk C.~S., White S.~D.~M., Springel V., Stadel J., Quinn T.~R., 2004, MNRAS, 349, 1039

\bibitem[Oh et~al.(2015)]{Oh2015}
Oh S.-H., et~al., 2015, AJ, 149, 180

\bibitem[Reines, Greene \& Geha(2013)]{ReinesGreeneGeha2013}
Reines A.~E., Greene J.~E., Geha M., 2013, ApJ, 775, 116

\bibitem[Schombert, McGaugh \& Lelli(2019)]{Schombert2019}
Schombert J.~M., McGaugh S.~S., Lelli F., 2019, MNRAS, 483, 1496

\bibitem[Schombert, McGaugh \& Lelli(2022)]{Schombert2022}
Schombert J.~M., McGaugh S.~S., Lelli F., 2022, AJ, 163, 154

\bibitem[Sellwood(2014)]{Sellwood2014}
Sellwood J.~A., 2014, Rev. Mod. Phys., 86, 1

\bibitem[Sellwood \& Wilkinson(1993)]{SellwoodWilkinson1993}
Sellwood J.~A., Wilkinson A., 1993, Rep. Prog. Phys., 56, 173

\end{thebibliography}